\subsection{Post-quantum security}
\label{sec:pq}

Every component of the shard and recursion arguments except the elliptic-curve multiset hash is hash-based: the Merkle commitments under $\Hash$, the \WHIR{} (or BaseFold) proximity test, the sumchecks and \LogUp{}. The Fiat--Shamir compilation of such interactive oracle proofs remains sound against quantum adversaries in the quantum random-oracle model~\citep{cms19}, with weaker concrete bounds: a quantum collision search on the 248-bit digest of the measured instantiation costs about $2^{83}$ rather than $2^{124}$ evaluations~\citep{bht98}, so the levels of \S\ref{sec:iopp-soundness} are classical. The digest is the exception. Shor's algorithm computes discrete logarithms in $E(\mathbb F_{p^7})$~\citep{shor97}, which reduces Assumption~\ref{ass:digest} to finding a short integer relation among known scalars, and at the interaction counts of a block such a relation should be assumed easy to find. Against a quantum adversary the options of \S\ref{sec:air-global} therefore differ: the global challenge and Merkle roots rest on hashes, and a lattice multiset hash on a problem not known to be easy for quantum adversaries, while the elliptic-curve multiset hash does not survive. Table~\ref{tab:pq} gives their costs; this paper claims post-quantum security for no configuration.

Of the two hash-based options, Merkle roots perform best: on one GPU, where the card is the bottleneck, their smaller committed area makes them faster; on four GPUs their proving work is smaller but arrives in execution order, and the cards idle while the last leaves and compositions of a block drain, a tail the elliptic-curve multiset hash fills with its independent precompile and memory shards.

\begin{remark}[Real-time quantum adversaries]
\label{rem:realtime}
Proving a block in time close to linear in its trace does not by itself protect the elliptic-curve digest. The lift applies no salt: a point is a fixed function of its message, so a quantum adversary can compute a vanishing combination once, offline, and reuse it in a false proof for any later block, and the speed of honest proving does not bound that attack. A time-bounded claim holds under two further conditions, neither of which the configuration of this paper meets. (i) \emph{Salting}: the lift mixes a per-block value $\sigma$ that is unknown before the block is fixed, such as the block hash or a challenge drawn after the shards' trace commitments, constrained in the chip and checked by the verifier, so that a relation among one block's points is useless for another. (ii) \emph{Timeliness}: a proof for block $b$ has value only if it is delivered within a window $\omega$ after $b$ appears, as in real-time proving where a late proof is not consumed. Under (i) and (ii) an adversary must find its relation for block $b$'s salted points within $\omega$, while the honest proof, and its verification, completes in about the honest proving time; if solving the discrete logarithms such a relation needs takes longer than $\omega$, the forgery arrives after the proof has lost its value, and fast proving lets $\omega$ shrink toward the honest proving time, 10 to 31\,s per block on four GPUs for the twelve blocks of Table~\ref{tab:pq}. This is a threat-model argument, not a theorem. It needs (i), which changes the digest and the verifying key; a lower bound on quantum discrete-logarithm time, which we do not have; and a deployment rule that enforces $\omega$. It gives no protection after $\omega$, and none for proofs that keep their value later, such as archived or bridged proofs.
\end{remark}
